\documentclass[11pt]{article}
\usepackage[margin=1in]{geometry}
\usepackage{amsmath,amssymb,amsthm}
\usepackage{graphicx}
\usepackage{hyperref}

\newtheorem{theorem}{Theorem}
\newtheorem{remark}{Remark}

\newcommand{\R}{\mathbb R}
\newcommand{\norm}[1]{\lVert #1\rVert}
\newcommand{\supp}{\operatorname{supp}}
\newcommand{\Sym}{\operatorname{Sym}}

\title{A Skew-Tensor Obstruction to Removing Symmetry in a Bourgain--Brezis Estimate}
\author{Run \texttt{fa\_banach\_001}, lane 8}
\date{July 18, 2026}

\begin{document}
\sloppy
\maketitle

\section*{Source Problem}

The source is Jean Van Schaftingen, \emph{Limiting Bourgain-Brezis estimates
for systems: theme and variations}, arXiv:1311.6624 \cite{VanSchaftingen}.
In Section 3.5, after stating the higher-order endpoint estimate on nilpotent
homogeneous groups for symmetric horizontal \(k\)-tensors, the paper says:

\begin{center}
\includegraphics[width=\textwidth]{figures/open_problem_crop.png}
\end{center}

\noindent
This packet answers the literal symmetry-removal question negatively: the
displayed estimate cannot hold if \(\varphi\) is allowed to be an arbitrary
horizontal \(k\)-tensor.  The obstruction already occurs for \(G=\R^2\), viewed
as a step-one stratified group, and \(k=2\).

\section*{Result}

\begin{theorem}\label{thm:main}
The higher-order estimate from Section 3.5 of the source paper becomes false
if the test tensor \(\varphi\in C_c^\infty(G;\Sym^k(T_bG))\) is replaced by an
arbitrary tensor \(\varphi\in C_c^\infty(G;\otimes^k T_bG)\).  In fact, it is
false already for the Euclidean group \(G=\R^2\), homogeneous dimension
\(Q=2\), and \(k=2\).
\end{theorem}

\section*{Proof Intuition}

The source constraint pairs \(f\) against second horizontal derivatives of
scalar functions.  In Euclidean space these second derivatives are Hessians,
hence symmetric bilinear forms.  Therefore a purely antisymmetric tensor field
automatically satisfies the constraint.  If the test tensor were also allowed
to be antisymmetric, the estimate would have to control the value of a
critical Sobolev function at a point by its \(L^2\) gradient norm.  Logarithmic
cutoffs show that such point control is impossible.

\section*{Proof}

\begin{proof}
Work on \(G=\R^2\), whose horizontal bundle is the whole tangent bundle and
whose homogeneous dimension is \(Q=2\).  Let
\[
  A=e_1\otimes e_2-e_2\otimes e_1
\]
be the standard nonzero antisymmetric 2-tensor.

Fix a nonnegative \(\rho\in C_c^\infty(B(0,1))\) with
\(\int_{\R^2}\rho=1\).  For \(m\ge1\), set \(r_m=e^{-2m}\) and
\[
  \rho_m(x)=r_m^{-2}\rho(x/r_m),\qquad f_m(x)=\rho_m(x)A .
\]
Then \(f_m\in C_c^\infty(\R^2;\otimes^2\R^2)\) and
\[
  \norm{f_m}_{L^1}=|A|\int_{\R^2}\rho_m=|A|.
\]
Moreover \(f_m\) satisfies the annihilation hypothesis from the higher-order
estimate.  Indeed, for every \(\psi\in C_c^\infty(\R^2)\), the Hessian
\(D^2\psi(x)\) is symmetric, and hence
\[
  \langle D^2\psi(x),A\rangle=0
  \qquad\text{for every }x\in\R^2.
\]
Therefore
\[
  \int_{\R^2}\langle D^2\psi,f_m\rangle=0.
\]

Assume, toward a contradiction, that the source estimate remained true with
arbitrary 2-tensor test fields.  Then there would be a constant \(C\) such
that, for all such \(f\) satisfying the above annihilation condition and all
\(\Phi\in C_c^\infty(\R^2;\otimes^2\R^2)\),
\[
  \left|\int_{\R^2}\langle \Phi,f\rangle\right|
  \le C\norm{f}_{L^1(\R^2)}\norm{\nabla\Phi}_{L^2(\R^2)}.
  \tag{1}
\]

Choose smooth radial cutoffs \(h_m\in C_c^\infty(B(0,2))\) with
\[
  h_m=1\quad\text{on }B(0,e^{-m})
\]
and
\[
  \norm{\nabla h_m}_{L^2(\R^2)}\le C_0m^{-1/2}.
\]
For example, smooth a standard logarithmic cutoff that equals
\(1\) for \(r\le e^{-m}\), equals \(\log(1/r)/m\) for
\(e^{-m}<r<1\), and equals \(0\) for \(r\ge1\).  The estimate follows from
\[
  \int_{e^{-m}<|x|<1}\frac{1}{m^2|x|^2}\,dx
  =\frac{2\pi}{m}.
\]

Now set \(\Phi_m=h_mA\).  Since \(r_m=e^{-2m}<e^{-m}\), we have
\(h_m=1\) on \(\supp\rho_m\).  Hence
\[
  \left|\int_{\R^2}\langle \Phi_m,f_m\rangle\right|
  =
  \left|\int_{\R^2} h_m(x)\rho_m(x)\langle A,A\rangle\,dx\right|
  =
  |A|^2.
\]
On the other hand,
\[
  \norm{\nabla\Phi_m}_{L^2}=|A|\norm{\nabla h_m}_{L^2}
  \le |A|C_0m^{-1/2}.
\]
Putting \(f=f_m\) and \(\Phi=\Phi_m\) in (1) gives
\[
  |A|^2
  \le C |A|^2 C_0m^{-1/2}
\]
for every \(m\), which is impossible as \(m\to\infty\).  Thus the estimate
cannot hold without the symmetry restriction on \(\Phi\).
\end{proof}

\section*{Limitations}

This is a counterexample to the naive or literal removal of the symmetry
restriction on the test tensor.  It does not rule out repaired variants that
add a separate condition annihilating or controlling antisymmetric components.
Nor does it settle the Besov and Sobolev--Lorentz open problems stated earlier
in the same source paper.

\section*{Novelty and Verification Notes}

The local registry, solution index, attempt index, and proof-gap index were
searched for arXiv:1311.6624 and for the phrases ``symmetry assumption'',
``symmetric tensor'', ``higher-order group estimate'', ``Chanillo Van
Schaftingen'', and ``antisymmetric''.  No local duplicate was found; the only
``antisymmetric'' hits were unrelated martingale and operator-space packets.
Web searches on 2026-07-18 for the exact source sentence and close variants
found the source paper or bibliographic mirrors, but no explicit counterexample
to the literal symmetry-removal extension.

For human review, the main interpretation issue is whether the source's
question means precisely replacing \(\Sym^k(T_bG)\) by the full tensor bundle
in the displayed estimate.  Under that reading, the proof above is a direct
counterexample.

\begin{thebibliography}{9}

\bibitem{VanSchaftingen}
Jean Van Schaftingen,
\emph{Limiting Bourgain-Brezis estimates for systems: theme and variations},
arXiv:1311.6624, 2013; J. Fixed Point Theory Appl. 15 (2014), no. 2,
273--297.

\end{thebibliography}

\end{document}
